\section{Casos prácticos y diagnóstico de problemas}
\subsection{Plantilla de Self-Feeding Chatbot}
Weston vuelve a revisar la práctica de Self-Feeding Chatbot: después de que el modelo dialoga con el usuario, convierte el draft inicial en una verified response, y luego incorpora las respuestas de alta calidad al conjunto de entrenamiento para apoyar el self-rewarding (SRT 774-884). Este paso enfatiza ``rewarding itself following instructions''; el modelo trata la verified response como otro tipo de label y la refuerza en future iterations.

\begin{table}[H]
\centering
\begin{tabular}{p{4cm}p{10cm}}
\toprule
Etapa & Descripción \\
\midrule
User Dialogue & El modelo interactúa con el usuario y registra la entrada, la salida y el contexto, lo que forma una señal de supervisión inmediata. \\
Model Draft & El modelo genera un draft y produce un Chain-of-Thought, que sirve como base del juicio. \\
Self-Verification & El modelo vuelve a revisar el draft, genera una verified response y enfatiza la consistencia lógica. \\
Verified Response & Esta respuesta se trata como una muestra de alta calidad para que el reward model aprenda de ella. \\
Training Update & La verified response se incorpora al conjunto de entrenamiento para seguir optimizando el reward / meta judge. \\
\bottomrule
\end{tabular}
\caption{Etapas clave y flujo de datos en la plantilla de Self-Feeding.}
\end{table}

\subsection{Patrones de error y alertas}
En las iteraciones de self-rewarding son comunes tres tipos de problemas: 1) overthinking del modelo, es decir, entre más piensa, más se aleja de la respuesta; 2) hallucination combinada con exceso de confianza; 3) meta judge drift, por ejemplo cuando un calibration drift interno del judge hace que incluso una verified response se juzgue erróneamente como mala. Weston también menciona ``we can use judge to cross check itself and see that it was wrong'' (SRT 1188-1190), lo que indica que necesitamos mantener un check-list confiable para detectar anomalías.

\begin{table}[H]
\centering
\begin{tabular}{p{4cm}p{5cm}p{5cm}}
\toprule
Patrón de error & Manifestación típica & Estrategia de respuesta \\
\midrule
Overthinking & El Chain-of-Thought se alarga gradualmente sin que la respuesta mejore & Limitar el reasoning depth e introducir un meta judge para corregir el rumbo \\
Hallucination + Confidence & El modelo se muestra seguro pero la respuesta se desvincula de los hechos & Usar múltiples evaluator y aumentar la autoverificación \\
Meta judge drift & Incluso una verified response se juzga como mala & Recalibrar periódicamente los datos del SelfT judge \\
\bottomrule
\end{tabular}
\caption{Alertas clave que deben monitorearse en la práctica de Self-Rewarding.}
\end{table}

\subsection{Resumen de la sección}
Self-Feeding Chatbot ofrece una plantilla de verificación en línea, y la tabla de patrones de error nos recuerda que es necesario recalibrar mediante el meta judge, limitar el reasoning depth y reforzar el monitoreo de hallucination para mantener el self-improvement en la dirección correcta.

